%!TeX root = ./../MusterAbschlussarbeit.tex

%##########################################################
% Inhalt
%##########################################################
\pagenumbering{arabic}
\chapter{Einleitung}

Der Leser Ihrer Arbeit sollte die Einleitung (auch als \enquote{Hintergrund} oder \enquote{Problemstellung} bezeichnet) als eine Art Rundblick erleben, 
der die Landschaft darstellt, mit der der Autor konfrontiert wurde, als die Arbeit entstand. Wann und wie hat sich dieses Forschungsthema erstmals angeboten? 
Was hat deutlich gemacht, dass es interessante Fragen gibt, die beantwortet werden müssen? Welche Handlungen und Zusammenhänge waren damals schon erkennbar, und welche Schlüsselinformationen fehlten? 

In der Regel wird in der Einleitung auch ein methodischer Ansatz zur Lösung des Problems vorgeschlagen, 
der möglicherweise durch konkrete Beispiele für die Anwendung dieses Ansatzes unter anderen Umständen veranschaulicht wird. 
Achten Sie beim Schreiben darauf, dass Sie klar zwischen Ihrer Arbeit und anderen Beiträgen unterscheiden müssen.

Bei umfangreichem Material, das Sie hier behandeln möchten, könnte es sich lohnen, die Einleitung in Unterabschnitte zu unterteilen

Kein anderer Abschnitt in der Arbeit hängt so sehr von der Literaturrecherche ab wie die Einleitung \cite{aoscw2004}. 
Bevor Sie mit der Arbeit an Ihrem Projekt begonnen haben, haben Sie sicher einige wichtige Hintergrundartikel gelesen, 
die Ihnen Ihr Betreuer schon empfohlen hat, aber Sie haben sie wahrscheinlich in letzter Zeit nicht mehr gelesen. 
Jetzt müssen Sie diese nicht nur erneut lesen, sondern sich auch eingehend mit den darin enthaltenen Verweisen befassen, 
um sich mit allen wichtigen Informationsquellen vertraut zu machen. Dies wird Sie natürlich auch in die Lage versetzen, 
frühere Arbeiten kritisch zu bewerten

\newpage

\section{Ein Beispielabschnitt mit diversen Dingen und Erläuterungen}
Zunächst ein Beispiel für Tabellen. Wichtig ist, dass die Tabellenbeschriftung a) da ist und b) \underline{über} der Tabelle steht. Die Lesenden wollen wissen was sie sehen, bevor sie in die Tabelle eintauchen.
Die hier gezeigte Tabelle ist ohne zusätzliche Pakete verwendbar. Alternativ wäre \texttt{longtable} einzubinden und zu verwenden.

\begin{table}[!h]
    \caption[Tabellenbeschriftung]{Tabellenbeschriftung \cite{Tane2014}}
    \centering
    \begin{tabular}{|c|c|c|}
        \hline
        ID & Kapitel & Seiten \\\hline
        0 & 1 & 5 \\\hline
    \end{tabular}
    \label{tab:Beispieltabelle}
\end{table}

Darüber hinaus drei Beispiele für Programm-Code. Im Ordner \glqq Programmcode\grqq{} sind drei Beispiele für C-, Python- und \ac{html}-Code, welche nachfolgend eingebunden ist.
Es empfiehlt sich ganze Dateien einzubinden, statt den Quellcode hier abzutippen.\\
Der Style ist fest definiert und muss nicht abgeändert werden. In aller Regel werden alle stilistischen Anliegen erfüllt.\\
Wie Sie im \LaTeX{} Quellcode sehen, ist die Caption direkt in den Aufruf reinzuschreiben. Gleiches gilt für ein Label, welches ich hier \ref{lst:C-Code-Beispiel}, hier \ref{lst:Python-Code-Beispiel} und hier \ref{lst:HTML-Code-Beispiel} referenziere.\\
Die Angabe der Sprache dient dem Zweck, dass sprachspezifische Keywords o.ä. geladen und erkannt werden, damit die Farbgebung passt.

\lstinputlisting[caption=C-Code Beispiel,label={lst:C-Code-Beispiel},language=C]{Programmcode/test.c}
\lstinputlisting[caption=Python-Code Beispiel,label={lst:Python-Code-Beispiel},language=Python]{Programmcode/test.py}

\textbf{Es ist dabei auch darauf zu achten, dass Quelltexte wenn möglich natürlich nicht über das Seitenende gehen sollten, wie z.B. beim Python-Code Bsp. Setzen Sie daher auch manuell
Seitenumbrüche.} Quelltexte welche eine Seite übersteigen, sind ohnehin eher als Anhang zu betrachten $\rightarrow$ siehe dazu Beschreibung im Anhang.

\lstinputlisting[caption=HTML-Code Beispiel,label={lst:HTML-Code-Beispiel},language=HTML]{Programmcode/test.html}

Nachfolgend noch ein einfaches Beispiel ein Bild einzubinden. Da es sich um Bild\textit{unterschriften} handelt, gehört diese somit \underline{unter} die Abbildung, anders als bei Tabellen\textit{überschriften}.\\
Im \LaTeX-Quelltext sieht man die Optionen für das Einbinden des Bildes. Neben \texttt{scale} wäre auch \texttt{width} möglich mit dem Parameter \texttt{linewidth} oder \texttt{textwidth}.\\
Achten Sie vor allem auf die Lesbarkeit der auf dem Bild befindlichen Informationen, unter der stetigen Annahme, dass es sich um ein ausgedrucktes Dokument handelt. Dort gibt es keinen Zoom.

Empfehlenswert wäre soweit möglich mit skalierbaren Vektorgrafiken zu arbeiten. Allerdings müssten Sie diese vor dem Aufruf von LaTeX in PDFs wandeln (inkscape IM-logo.svg -o IM-logo.pdf) oder Inkscape installieren und LaTeX mit --shell-escape aufrufen.

\begin{figure}[!t]
    \centering
    \includegraphics[scale=0.6]{Bilder/logos/IM-logo}
    \caption{Das Logo von Informatik und Medien }
\end{figure}

Bedenken Sie bei allen Beschriftungen, egal ob Abbildungen, Tabellen oder Quelltexte, dass Sie, sofern diese nicht Ihrer eigenen Schaffenskraft entsprangen,
diese referenzieren. Im Beispiel der Tabelle \ref{tab:Beispieltabelle} sieht man, dass sich der Beschreibungstext direkt an der Tabelle von dem im Tabellenverzeichnis unterscheidet,
genau um den Punkt der Referenz.

\section{Methodik}

Im Methodik-Abschnitt (eine eigene Section) einer Arbeit stellen Sie die Methoden dar, die Sie verwendet haben oder vorhaben zu verwenden, um die Fragestellung zu beantworten. 
In diesem Abschnitt erläutern Sie, welche Methoden umgesetzt wurden (oder auch angedacht), um die Hypothesen zu testen und die Fallstudie durchzuführen. 
Welche Methoden für Ihre Arbeit geeignet sind, ist auch ein wesentlicher Punkt Ihres angestrebten Titels und somit ein zu bewertender Teil dieser Arbeit. 

\underline{Aber \textbf{mindestens} folgende Punkte}\\
Inhalt des ersten Kapitels (im Allgemeinen):
\begin{itemize}
    \item Thematische Einführung, umreißen des Themas \dots wo befinden wir uns?
    \item Herausstellen des Problems und warum es eines ist was gelöst werden sollte.
    \item Herausstellen der tatsächlichen Ziel-/Fragestellung, welche bearbeitet werden wird.
    \item Eine Abgrenzung was diese Arbeit ist und was sie nicht ist.
    \item Stand der Forschung \dots Was existiert bereits? Wie gut passt das auf das Problem?
    \begin{itemize}
        \item Referenzen mittels cite: \cite[S.~111]{jsch2011} \cite[S.~27f]{Tane2014}
    \end{itemize}
    \item Eine abgeleitete Methodik $\rightarrow$ basierend auf der Literatur und Ihrem Wissen: wie haben Sie nun vor Ihr Problem zu lösen? Vllt durch \cite{9429985}?
\end{itemize}



